% Zdroj: PDF priklady-podzim-2025.pdf, fyzická str. 1-2. Značka: společné okruhy. Zařazeno: I4.
% Pozn.: v originálu jsou vícebajtová čísla v BIG-endian pořadí (ne little-endian).
\section{Ukládání dat v binárním souboru}
\szzotazka{podzim 2025}{2}{Ukládání dat v~binárním souboru (hex dump, reprezentace dat)}

\noindent\textbf{Zadání.}\par
Aplikace ukládá svá data do binárního souboru, což je ilustrováno následujícím (částečným)
hexadecimálním výpisem:

\begin{lstlisting}
0000:    00 06 48 68     53 53 4C 4C     00 22 03 E8     FC 18 00 05     ..Hh SSLL .".. ....
0010:    48 65 6C 6C     6F 00 05 57     6F 72 6C 64     00 00 00 00     Hell o..W orld ....
0020:    00 00 00 2C     00 00 00 00     00 00 12 34     00 88 01 F4     ..., .... ...4 ....
\end{lstlisting}

Logicky soubor obsahuje \textit{uzly}. Každý uzel obsahuje stejnou sadu \textit{atributů}.
Typy atributů jsou určeny \textit{hlavičkou} souboru umístěnou na jeho začátku, která se
skládá ze dvou bajtů udávajících počet atributů (v~našem příkladu 6), následovaných typy
atributů, kódovanými jedním bajtem pro každý atribut, odpovídající tomuto výčtu:

\begin{lstlisting}[language=C++]
enum AttrType { TUInt16 = 0x48, TLink = 0x4C, TString = 0x53, TSInt16 = 0x68};
\end{lstlisting}

Vícebajtová čísla uložená v~souboru jsou vždy v~big-endian pořadí. Bezprostředně za
hlavičkou se nachází \textit{kořenový uzel}. Každý uzel začíná dvěma bajty obsahujícími
délku dat uzlu v~bajtech (v~kořenovém uzlu našeho příkladu 34). Data uzlu obsahují hodnoty
atributů, kódované podle typů atributů deklarovaných v~hlavičce. Pro typy \texttt{TUInt16}
a~\texttt{TSInt16} je atribut uložen jako dva bajty obsahující 16bitové celé číslo bez
znaménka nebo se znaménkem (dvojkový doplněk). Atributy typu \texttt{TString} jsou ASCII
řetězce proměnné délky kódované jako dva bajty určující počet znaků, následované znaky,
každý uložený v~jednom bajtu. Pro \texttt{TLink} je uloženo 8 bajtů obsahujících pozici
jiného uzlu v~souboru. Částečný výpis ukazuje, že kořenový uzel nese hodnoty
$\{1000, -1000, \texttt{\textquotedbl Hello\textquotedbl}, \texttt{\textquotedbl World\textquotedbl}, \texttt{0x2C}, \texttt{0x1234}\}$.

Binární soubor je reprezentován třídou \texttt{BinaryFile}, která poskytuje následující
metodu, jež načte \texttt{len} bajtů do bufferu \texttt{buf}, začínajícího na absolutní
pozici \texttt{filepos} v~souboru (znaménkové typy se používají kvůli kompatibilitě
s~Javou):

\begin{lstlisting}[language=C++]
void readBytes (long filepos, byte[] buf, int len);                  // Java, C#
void readBytes (std::int64_t filepos, unsigned char* buf, int len);  // C++
\end{lstlisting}

\begin{enumerate}
\item Napište deklaraci třídy \texttt{Reader}, včetně jejích privátních dat
  a~implementace následujících metod: Konstruktor obdrží otevřený objekt
  \texttt{BinaryFile} jako parametr; okamžitě načte hlavičku. Metoda \texttt{attributes}
  vrací počet atributů. Metoda \texttt{getType} vrací typ $i$-tého atributu, číslováno od
  nuly. Metoda \texttt{getRoot} vrací pozici kořenového uzlu. Metoda \texttt{readNode}
  vrátí nový objekt \texttt{Node} reprezentující uzel uložený na pozici \texttt{filepos}.
  Objekt je vracen odkazem v~Java/C\#, ale hodnotou v~C++. Data uzlu se načítají ze
  souboru pouze touto funkcí.
\item Deklarujte a~implementujte privátní datové prvky a~konstruktor třídy \texttt{Node}.
  Třída má ukládat data uzlu jako pole bajtů, tj. tak, jak jsou uložena v~souboru. Navíc
  má obsahovat strukturu, která umožní lokalizovat libovolný atribut v~konstantním čase.
  Konverze bajtů na celá čísla nebo řetězce se provádí pouze při volání odpovídajících
  metod \texttt{get...} třídy \texttt{Node}.
\item Napište implementaci následujících metod třídy \texttt{Node} pro získání hodnoty
  celočíselných atributů se znaménkem a~bez znaménka (16bit), použijte přitom pouze
  vestavěné operátory zvoleného jazyka:

\begin{lstlisting}[language=C++]
int getUInt16(int i);
int getSInt16(int i);
\end{lstlisting}

Každá metoda čte $i$-tý atribut; je odpovědností volajícího zajistit, že volání odpovídá
typu atributu. Typ \texttt{int} se považuje za dostatečně velký, aby obsahoval rozsah
$\langle -32768, 65535 \rangle$.
\end{enumerate}

\medskip
\noindent\textbf{Náčrt řešení.}\par
\textbf{(0)~Dekódování výpisu.} Hlavička: \texttt{00 06} $=6$ atributů, pak typy
\texttt{48 68 53 53 4C 4C} $=$ \texttt{TUInt16, TSInt16, TString, TString, TLink, TLink}.
Kořenový uzel začíná na offsetu~8 (za 2~B počtu a~6~B typů), jeho délka je \texttt{00 22} $=34$~B.
Data (big-endian): \texttt{03 E8}$=1000$; \texttt{FC 18}$=\texttt{0xFC18}-2^{16}=-1000$;
délka~5 + \texttt{Hello}; délka~5 + \texttt{World}; 8~B \texttt{...00 2C}$=\texttt{0x2C}$;
8~B \texttt{...12 34}$=\texttt{0x1234}$. Spotřeba $2{+}2{+}7{+}7{+}8{+}8=34$~B sedí.

\textbf{(1)~Reader.} Konstruktor přečte počet atributů a~typy; \texttt{getRoot} vrací offset hned za
hlavičkou. \texttt{readNode} přečte 2~B délky a~pak data uzlu.
\begin{lstlisting}[language=C++]
struct Reader {
  BinaryFile& f;  int nattr;  std::vector<unsigned char> types;
  Reader(BinaryFile& bf) : f(bf) {
    unsigned char h[2];  f.readBytes(0, h, 2);
    nattr = (h[0] << 8) | h[1];                  // big-endian
    types.resize(nattr);  f.readBytes(2, types.data(), nattr);
  }
  int  attributes()   { return nattr; }
  int  getType(int i) { return types[i]; }
  long getRoot()      { return 2 + nattr; }      // hned za hlavickou
  Node readNode(long filepos) {
    unsigned char len2[2];  f.readBytes(filepos, len2, 2);
    int len = (len2[0] << 8) | len2[1];
    std::vector<unsigned char> data(len);
    f.readBytes(filepos + 2, data.data(), len);
    return Node(types, std::move(data));
  }
};
\end{lstlisting}

\textbf{(2)--(3)~Node.} Data drží jako pole bajtů; v~konstruktoru předpočítá offsety atributů
(kvůli proměnné délce řetězců), takže $i$-tý atribut najde v~$O(1)$. \texttt{get} metody skládají
big-endian.
\begin{lstlisting}[language=C++]
struct Node {
  std::vector<unsigned char> data;  std::vector<int> off;   // off[i] = pozice atributu i
  Node(const std::vector<unsigned char>& types, std::vector<unsigned char> d)
      : data(std::move(d)) {
    int p = 0;
    for (unsigned char t : types) {
      off.push_back(p);
      if (t == 0x48 || t == 0x68) p += 2;                   // TUInt16 / TSInt16
      else if (t == 0x4C)         p += 8;                   // TLink
      else if (t == 0x53)         p += 2 + ((data[p]<<8)|data[p+1]);  // TString
    }
  }
  int getUInt16(int i) { int p=off[i]; return (data[p]<<8) | data[p+1]; }
  int getSInt16(int i) { int v=getUInt16(i); return v>=0x8000 ? v-0x10000 : v; }
};
\end{lstlisting}
\textit{(Náčrt doplněn k~výkladu okruhu; otázka nemá v~PDF řešení, dekódování i~kód ověřeny skriptem \texttt{verify\_szz.py}.)}
